% TOP_PROBLEM: {"id": 3246, "title": "Colouring the square of a planar graph", "category_id": 3, "category": {"id": 3, "name": "graph_theory", "display_name": "Graph Theory", "description": "Problems involving graphs, networks, and their properties.", "slug": "graph-theory", "order_index": 3, "created_at": "2026-07-31T15:26:25.671Z"}, "status": "open", "problem_number": "OPG-47470", "set_id": 12}
% FIRST_POSTED: 2026-10-01
% ENTRY_KIND: statement_only
% UnsolvedMath ID 3246; source catalogue code OPG-47470
% !TeX program = lualatex
% Definition and sources only; this file is not an LLM solution attempt.
\documentclass[11pt,a4paper]{article}
\usepackage[margin=25mm]{geometry}
\usepackage{fontspec}
\setmainfont{lmroman10-regular.otf}[BoldFont=lmroman10-bold.otf,
 ItalicFont=lmroman10-italic.otf,BoldItalicFont=lmroman10-bolditalic.otf]
\usepackage{amsmath,amssymb,mathtools}
\usepackage{unicode-math}
\setmathfont{latinmodern-math.otf}
\AtBeginDocument{\let\setminus\smallsetminus}
\usepackage{xurl,hyperref}
\hypersetup{colorlinks=true,urlcolor=blue}
\setcounter{secnumdepth}{0}
\setlength{\parindent}{0pt}
\setlength{\parskip}{5pt}
\setlength{\emergencystretch}{3em}
\begin{document}
\section*{Colouring the square of a planar graph}\label{3246}
{\small UnsolvedMath ID 3246 · OPG-47470 · Graph Theory · OpenGarden}
\subsection{Definitions and mathematical statement}
For a finite simple undirected graph $G$, the distance between two vertices
in the same component is the minimum number of edges in a path joining
them; vertices in different components have infinite distance. The square
$G^2$ has vertex set $V(G)$ and an edge between distinct vertices exactly
when their distance in $G$ is at most two. Write $\Delta=\Delta(G)$ for
the maximum degree of the original graph, and $\chi(H)$ for the minimum
number of colours in a proper colouring of a graph $H$.

The recorded conjecture asks whether every finite simple planar graph
with $\Delta\ge3$ satisfies
\[
 \chi(G^2)\le
 \begin{cases}
  7,&\Delta=3,\\
  \Delta+5,&4\le\Delta\le7,\\
  \lfloor3\Delta/2\rfloor+1,&\Delta\ge8.
 \end{cases}
\]
Thus two vertices must receive different colours if they are adjacent
or have a common neighbour in $G$. The square itself need not be planar;
planarity is assumed only of $G$. The floor symbol denotes rounding down
to an integer.

The parameter is ordinary chromatic number: a common palette is chosen
for the whole graph. The source also discusses a stronger list-colouring
variant, but arbitrary lists are not an additional requirement in the
three displayed bounds. The statement includes all finite orders,
disconnected graphs, and graphs with vertices of degree below $\Delta$;
no regularity is assumed. Degrees zero, one, and two are outside the
listed regimes and do not change the three target assertions.
\subsection{Short English statement}
Colour vertices of a planar graph so that vertices at distance one or two differ, using the stated degree-dependent number of colours.
\subsection{Sources}
\begin{itemize}
\item[S1] ULAM.AI and the UnsolvedMath Contributors, supplied problems.json, record 3246 (OPG-47470), OpenGarden collection. Catalogue identity and source transcription; CC BY 4.0.
\url{https://huggingface.co/datasets/ulamai/UnsolvedMath}.
\item[S2] Open Problem Garden, Colouring the square of a planar graph. Original problem and discussion; the mathematical formulation below is expanded from the supplied source record.
\url{http://www.openproblemgarden.org/op/colouring_the_square_of_a_planar_graph}.
\end{itemize}
\end{document}
